% 06_simulation.tex - ROS Simulation Module
% AlohaMini Codebase Documentation

\section{ROS Simulation}

% ============================================
% SLIDE: Module Overview
% ============================================
\begin{frame}{simulation/ Module Overview}
\begin{columns}
\begin{column}{0.4\textwidth}
    \textbf{Purpose:} ROS visualization \& physics

    \vspace{0.2cm}

    \textbf{Structure:}
    \begin{itemize}
        \item \filepath{src/Aloha/} - ROS package
        \item \filepath{urdf/} - Robot model
        \item \filepath{meshes/} - STL files
        \item \filepath{launch/} - Launch files
    \end{itemize}

    \vspace{0.2cm}

    \textbf{Support:}
    \begin{itemize}
        \item ROS1 Noetic / ROS2 Humble
        \item RViz / Gazebo
    \end{itemize}
\end{column}
\begin{column}{0.6\textwidth}
    \centering
    \begin{tikzpicture}[scale=0.7, transform shape]
        % Package structure
        \node[module, text width=3.2cm] (pkg) at (0,0) {\textbf{Aloha Package}};

        % Components
        \node[submodule, text width=2cm, below left=1cm and 0.3cm of pkg] (urdf) {urdf/};
        \node[submodule, text width=2cm, below=1cm of pkg] (launch) {launch/};
        \node[submodule, text width=2cm, below right=1cm and 0.3cm of pkg] (meshes) {meshes/};

        % Arrows
        \draw[arrow] (pkg) -- (urdf);
        \draw[arrow] (pkg) -- (launch);
        \draw[arrow] (pkg) -- (meshes);

        % Tools
        \node[class, text width=1.8cm, above left=0.6cm and 0cm of pkg] (rviz) {RViz};
        \node[class, text width=1.8cm, above right=0.6cm and 0cm of pkg] (gazebo) {Gazebo};

        \draw[flow] (pkg) -- (rviz);
        \draw[flow] (pkg) -- (gazebo);
    \end{tikzpicture}
\end{column}
\end{columns}
\end{frame}

% ============================================
% SLIDE: URDF Structure
% ============================================
\begin{frame}{URDF Robot Model}
\begin{columns}
\begin{column}{0.5\textwidth}
    \textbf{Link/Joint Structure:}
    \begin{itemize}
        \item \code{base\_link} - Robot base
        \item \code{wheel1-3} - Omniwheels
        \item \code{vertical\_move} - Lift
        \item \code{left\_joint1..6} - Left arm
        \item \code{right\_joint1..6} - Right arm
    \end{itemize}

    \vspace{0.2cm}

    \textbf{Total:} 17 joints (3 + 1 + 6 + 6 + gripper)
\end{column}
\begin{column}{0.5\textwidth}
    \centering
    \begin{tikzpicture}[scale=0.85, transform shape, level distance=1.4cm,
        level 1/.style={sibling distance=3.5cm},
        level 2/.style={sibling distance=2.2cm}]
        \node[module, text width=2cm] {base\_link}
            child {node[submodule, text width=1.7cm] {\scriptsize wheels (3)}}
            child {node[module, text width=2cm] {vertical\_link}
                child {node[submodule, text width=1.7cm] {\scriptsize left\_arm (6)}}
                child {node[submodule, text width=1.7cm] {\scriptsize right\_arm (6)}}
            };
    \end{tikzpicture}
\end{column}
\end{columns}
\end{frame}

% ============================================
% SLIDE: Joint Configuration
% ============================================
\begin{frame}[fragile]{Joint Names Configuration}
\begin{lstlisting}[basicstyle=\ttfamily\scriptsize]
# config/joint_names_Aloha.yaml
controller_joint_names: [
    '',              # Empty (index 0)
    'wheel1', 'wheel2', 'wheel3',         # Wheels
    'vertical_move',                       # Lift
    'left_joint1', ..., 'left_joint6',    # Left arm
    'right_joint1', ..., 'right_joint6',  # Right arm
]
\end{lstlisting}

\vspace{0.3cm}

\textbf{Joint Indexing:}
\begin{itemize}
    \item Index 0: Empty placeholder
    \item Index 1-3: Omniwheels
    \item Index 4: Lift mechanism
    \item Index 5-10: Left arm (6 DOF)
    \item Index 11-16: Right arm (6 DOF)
\end{itemize}
\end{frame}

% ============================================
% SLIDE: ROS1 Launch Files
% ============================================
\begin{frame}[fragile]{ROS1 Launch Files}
\begin{columns}
\begin{column}{0.5\textwidth}
\textbf{display.launch} (RViz):
\begin{lstlisting}
<launch>
  <param name="robot_description"
    textfile="$(find Aloha)/urdf/Aloha.urdf"/>

  <node name="joint_state_publisher_gui"
    pkg="joint_state_publisher_gui"
    type="joint_state_publisher_gui"/>

  <node name="robot_state_publisher"
    pkg="robot_state_publisher"
    type="robot_state_publisher"/>

  <node name="rviz"
    pkg="rviz"
    type="rviz"
    args="-d $(find Aloha)/urdf.rviz"/>
</launch>
\end{lstlisting}
\end{column}
\begin{column}{0.5\textwidth}
\textbf{gazebo.launch} (Physics):
\begin{lstlisting}
<launch>
  <include file="$(find gazebo_ros)/
    launch/empty_world.launch"/>

  <node name="tf_footprint_base"
    pkg="tf"
    type="static_transform_publisher"
    args="0 0 0 0 0 0 base_link
          base_footprint 40"/>

  <node name="spawn_model"
    pkg="gazebo_ros"
    type="spawn_model"
    args="-file $(find Aloha)/urdf/
          Aloha.urdf -urdf -model Aloha"
    output="screen"/>
</launch>
\end{lstlisting}
\end{column}
\end{columns}
\end{frame}

% ============================================
% SLIDE: ROS2 Launch File
% ============================================
\begin{frame}[fragile]{ROS2 Launch File (Python)}
\begin{lstlisting}
# display.launch.py
from ament_index_python.packages import get_package_share_directory
from launch import LaunchDescription
from launch_ros.actions import Node

def generate_launch_description():
    bringup_dir = get_package_share_directory('Aloha')
    urdf_path = os.path.join(bringup_dir, 'urdf', 'Aloha.urdf')

    with open(urdf_path, 'r') as f:
        robot_description = f.read()

    return LaunchDescription([
        Node(package='robot_state_publisher', executable='robot_state_publisher',
             parameters=[{'robot_description': robot_description}]),
        Node(package='joint_state_publisher', executable='joint_state_publisher'),
        Node(package='rviz2', executable='rviz2',
             arguments=['-d', rviz_config_path]),
    ])
\end{lstlisting}
\end{frame}

% ============================================
% SLIDE: Package Dependencies
% ============================================
\begin{frame}[fragile]{Package Dependencies (package.xml)}
\begin{lstlisting}
<?xml version="1.0"?>
<package format="3">
  <name>Aloha</name>
  <version>0.1.0</version>
  <description>AlohaMini Robot</description>

  <buildtool_depend>ament_cmake</buildtool_depend>

  <exec_depend>joint_state_publisher</exec_depend>
  <exec_depend>joint_state_publisher_gui</exec_depend>
  <exec_depend>robot_state_publisher</exec_depend>
  <exec_depend>rviz2</exec_depend>
  <exec_depend>xacro</exec_depend>

  <export>
    <build_type>ament_cmake</build_type>
  </export>
</package>
\end{lstlisting}
\end{frame}

% ============================================
% SLIDE: RViz Visualization
% ============================================
\begin{frame}{RViz Visualization}
\begin{columns}
\begin{column}{0.5\textwidth}
    \textbf{Features:}
    \begin{itemize}
        \item Robot model visualization
        \item TF frame tree display
        \item Joint state publisher GUI
        \item Interactive joint control
    \end{itemize}

    \vspace{0.3cm}

    \textbf{Usage (ROS1):}
    \begin{itemize}
        \item \code{roslaunch Aloha display.launch}
    \end{itemize}

    \textbf{Usage (ROS2):}
    \begin{itemize}
        \item \code{ros2 launch Aloha display.launch.py}
    \end{itemize}
\end{column}
\begin{column}{0.5\textwidth}
    \centering
    \begin{tikzpicture}[scale=0.8, transform shape, node distance=1cm]
        % RViz workflow
        \node[module, text width=3cm] (urdf) {URDF Model};
        \node[module, text width=3cm, below=of urdf] (rsp) {Robot State\\Publisher};
        \node[module, text width=3cm, below=of rsp] (jsp) {Joint State\\Publisher};
        \node[module, text width=3cm, below=of jsp] (rviz) {RViz};

        \draw[flow] (urdf) -- (rsp);
        \draw[flow] (rsp) -- (jsp);
        \draw[flow] (jsp) -- (rviz);

        % Topics
        \node[submodule, text width=2cm, right=0.5cm of rsp] (tf) {/tf};
        \node[submodule, text width=2cm, right=0.5cm of jsp] (js) {/joint\_states};

        \draw[flow, dashed] (rsp) -- (tf);
        \draw[flow, dashed] (jsp) -- (js);
        \draw[flow, dashed] (tf) -- (rviz);
        \draw[flow, dashed] (js) -- (rviz);
    \end{tikzpicture}
\end{column}
\end{columns}
\end{frame}

% ============================================
% SLIDE: Gazebo Simulation
% ============================================
\begin{frame}[shrink=3]{Gazebo Physics Simulation}
\begin{columns}
\begin{column}{0.5\textwidth}
    \textbf{Features:}
    \begin{itemize}
        \item Full physics simulation
        \item Collision detection
        \item Inertia-based dynamics
        \item World interaction
    \end{itemize}

    \vspace{0.3cm}

    \textbf{Launch Components:}
    \begin{enumerate}
        \item Empty world spawn
        \item TF footprint publisher
        \item Model spawner
        \item Fake joint calibration
    \end{enumerate}

    \vspace{0.3cm}

    \textbf{Usage:}
    \begin{itemize}
        \item \code{roslaunch Aloha gazebo.launch}
    \end{itemize}
\end{column}
\begin{column}{0.5\textwidth}
    \centering
    \begin{tikzpicture}[scale=0.9, transform shape, node distance=1.1cm]
        % Gazebo workflow
        \node[module, text width=3.2cm] (world) {Gazebo World};

        % Physics Engine
        \node[class, text width=2.6cm, right=0.6cm of world] (physics) {Physics\\Engine};
        \draw[flow, <->] (world) -- (physics);

        % Spawn Model
        \node[module, text width=3.2cm, below=of world] (spawn) {Spawn Model};
        \draw[flow] (spawn) -- (world);

        % URDF and Meshes
        \node[submodule, text width=2.2cm, below left=0.9cm and -0.3cm of spawn] (urdf) {URDF};
        \node[submodule, text width=2.2cm, below right=0.9cm and -0.3cm of spawn] (mesh) {Meshes};

        \draw[flow] (urdf) -- (spawn);
        \draw[flow] (mesh) -- (spawn);

        % Topics output
        \node[submodule, text width=2.8cm, below=2.8cm of world] (topics) {
            \scriptsize /gazebo/states\\
            \scriptsize /calibrated
        };
        \draw[flow, dashed] (world) -- node[left, font=\tiny] {publish} (topics);
    \end{tikzpicture}
\end{column}
\end{columns}
\end{frame}

% ============================================
% SLIDE: Directory Structure
% ============================================
\begin{frame}[fragile,shrink=3]{Directory Structure}
\begin{columns}
\begin{column}{0.5\textwidth}
\begin{lstlisting}[basicstyle=\ttfamily\tiny,numbers=none,frame=none]
simulation/src/Aloha/
+-- config/
|   +-- joint_names_Aloha.yaml
+-- launch/
|   +-- display.launch      # ROS1
|   +-- display.launch.py   # ROS2
|   +-- gazebo.launch
+-- meshes/
|   +-- base_link.STL
|   +-- wheel*.STL
|   +-- left_joint*.STL
|   +-- right_joint*.STL
+-- urdf/
|   +-- Aloha.urdf
+-- package.xml
\end{lstlisting}
\end{column}
\begin{column}{0.5\textwidth}
    \textbf{Mesh Files:}
    \begin{itemize}
        \item STL format (binary)
        \item SolidWorks export
        \item Collision \& visual
    \end{itemize}

    \vspace{0.2cm}

    \textbf{Build System:}
    \begin{itemize}
        \item ament\_cmake (ROS2)
        \item catkin (ROS1 compatible)
    \end{itemize}
\end{column}
\end{columns}
\end{frame}

% ============================================
% SLIDE: Module Summary
% ============================================
\begin{frame}{ROS Simulation Summary}
\begin{columns}
\begin{column}{0.5\textwidth}
    \textbf{Key Components:}
    \begin{itemize}
        \item URDF robot model
        \item STL mesh files
        \item ROS1/ROS2 launch files
        \item Joint configuration
    \end{itemize}

    \vspace{0.3cm}

    \textbf{Visualization:}
    \begin{itemize}
        \item RViz for visual inspection
        \item Joint state publisher GUI
        \item TF frame visualization
    \end{itemize}
\end{column}
\begin{column}{0.5\textwidth}
    \textbf{Physics Simulation:}
    \begin{itemize}
        \item Gazebo integration
        \item Inertia-based dynamics
        \item Collision meshes
    \end{itemize}

    \vspace{0.3cm}

    \textbf{Comparison with ManiSkill3:}
    \begin{table}
    \centering
    \footnotesize
    \begin{tabular}{lcc}
    \toprule
    \textbf{Feature} & \textbf{ROS} & \textbf{ManiSkill3} \\
    \midrule
    Physics & Gazebo & SAPIEN \\
    GPU Accel & No & Yes \\
    RL Support & Limited & Native \\
    Visualization & RViz & Viewer \\
    \bottomrule
    \end{tabular}
    \end{table}
\end{column}
\end{columns}
\end{frame}

